% !TeX program = xelatex
% !TeX encoding = UTF-8
\documentclass{article}
\usepackage[paperwidth=595.3pt,paperheight=841.9pt,left=85.05pt,right=42.5pt,top=26.4pt,bottom=14.2pt,nohead,nofoot]{geometry}
\usepackage{fontspec}
\setmainfont{Times New Roman}[Ligatures=TeXOff]
\usepackage[normalem]{ulem}
\usepackage{xcolor}
\usepackage{graphicx}

\pagestyle{empty}
\parskip=0pt \parindent=0pt \lineskip=0pt \normallineskip=0pt \topskip=0pt
\frenchspacing \sloppy \raggedbottom
\hyphenpenalty=10000 \hbadness=10000

\newdimen\wL \newdimen\wR \newdimen\wB \newdimen\wA \newdimen\wBd \newdimen\wX
\newcommand{\wcenter}{\def\walign{c}}
\newcommand{\wright}{\def\walign{r}}
\newcommand{\wjustify}{\def\walign{j}}
\newcommand{\wind}[3]{\wL=#1 \wR=#2 \parindent=#3 }
\newcommand{\wspace}[2]{\wB=#1 \wA=#2 }
\newcommand{\wrule}[1]{\wBd=#1 }
\newcommand{\wsz}[2]{\fontsize{#1}{#2}\selectfont
  \def\wextra{\dimexpr #2pt-1.15\dimexpr #1pt\relax\relax}%
  \def\wstrut{\vrule width 0pt height 0.93\dimexpr #1pt\relax depth \dimexpr #2pt-0.93\dimexpr #1pt\relax\relax}}
\newcommand{\wfix}{\spaceskip=\fontdimen2\font \if\walign j\advance\spaceskip 0pt plus \fontdimen3\font \fi}
\makeatletter
\newenvironment{wpar}[1][]{%
  \par \wL=0pt \wR=0pt \wB=0pt \wA=0pt \wBd=0pt \parindent=0pt \def\walign{l}#1%
  \vskip\wB
  \leftskip=\wL \rightskip=\wR \parfillskip=0pt plus 1fil
  \if\walign l\advance\rightskip 0pt plus 1fil \fi
  \if\walign c\advance\leftskip 0pt plus 1fil \advance\rightskip 0pt plus 1fil \parfillskip=0pt \fi
  \if\walign r\advance\leftskip 0pt plus 1fil \parfillskip=0pt \fi
  \wfix\@ifnextchar\wstretch{}{\leavevmode\wstrut\ignorespaces}%
}{%
  \ifhmode\unskip\wstrut\fi\par\nointerlineskip
  \ifdim\wBd>0pt \vskip\dimexpr1pt-\wextra\relax
    \moveright\wL\vbox{\hrule height\wBd width\dimexpr\hsize-\wL-\wR\relax}\nointerlineskip\vskip\wextra\fi
  \vskip\wA
}
\makeatother
\newcommand{\пусто}[1][]{\begin{wpar}[#1]\mbox{}\end{wpar}} %
\newcount\wN
\newcommand{\wspaces}[1]{\wN=#1 \loop\ifnum\wN>0 \hskip\spaceskip \advance\wN by -1 \repeat}
\renewcommand{\ULthickness}{0.05em}
\newcommand{\wul}[1]{{\ULdepth=0.16em \uline{#1}}}
\newcommand{\wulp}[1]{\expandafter\wul\expandafter{#1}}
\makeatletter
\newcommand{\wsmall}[1]{\fontsize{0.58\dimexpr\f@size pt\relax}{0pt}\selectfont\wfix #1}
\makeatother
\newcommand{\wsup}[1]{\raisebox{0.33em}{\wsmall{#1}}}
\newcommand{\wsub}[1]{\raisebox{-0.08em}{\wsmall{#1}}}
\newcommand{\wsize}[2]{{\fontsize{#1}{#1}\selectfont\wfix #2}}
\newcommand{\wcolor}[2]{\textcolor[HTML]{#1}{#2}}
\newcommand{\wfont}[2]{{\fontspec{#1}#2}}
\newcommand{\wdash}[2][\_]{\wN=#2 \wX=0pt \setbox0\hbox{\_}%
  \loop\ifnum\wN>0 \advance\wX\wd0
    \ifdim\wX>\dimexpr\hsize-\wL-\wR\relax \break \wX=\wd0 \fi
    #1\advance\wN by -1 \repeat}
\newcommand{\wstretch}[1]{%
  \setbox2\vbox{\leavevmode\wstrut #1\unskip\par \setbox0\lastbox \nointerlineskip
    \hbox to\hsize{\unhbox0 \unskip\unskip\unpenalty\hskip\rightskip}\global\wX=\prevdepth}%
  \unvbox2 \prevdepth=\wX \noindent}
\newdimen\wtabstop \wtabstop=35.4pt
\newcount\wT
\newcommand{\wtab}[1]{%
  \setbox0\hbox{#1}\wX=\dimexpr\wd0+\wL+\parindent\relax
  \ifdim\wX<\wL \wX=\dimexpr\wL-\wX\relax %
  \else \wT=\wX \divide\wT by \wtabstop \advance\wT by 1 \wX=\dimexpr\wT\wtabstop-\wX\relax \fi
  \unhbox0\hskip\wX\relax}
\newcommand{\wcap}[3][13.8pt]{\setbox0\hbox{#2}%
  \rlap{\raisebox{-#1}[0pt][0pt]{\hbox to\wd0{\hss#3\hss}}}\unhbox0}
\newcommand{\wfield}[1]{%
  \setbox0\hbox{#1}%
  \wcap{\hbox{\wul{\phantom{\copy0}}}}{\wsub{#1}}}
\newenvironment{wtable}[1]{\par\nointerlineskip\noindent\hbox to\linewidth\bgroup\hskip0pt\relax
  \setlength{\tabcolsep}{5.4pt}\renewcommand{\arraystretch}{0}\begin{tabular}{#1}}%
  {\end{tabular}\hss\egroup\par\nointerlineskip}
\newcommand{\wcell}[2]{\parbox[t]{\dimexpr #1-10.8pt\relax}{#2}}
\newcommand{\wrowheight}[1]{\vrule width 0pt depth #1}
\newcommand{\wlabel}[2]{\makebox[#1][l]{#2}}
\newsavebox{\wboxbox}
\newenvironment{wbox}[1]{\par\nointerlineskip\setlength{\fboxsep}{0pt}\setlength{\fboxrule}{0.5pt}%
  \noindent\hskip\dimexpr#1-4.5pt\relax
  \begin{lrbox}{\wboxbox}\hskip4pt\begin{minipage}{\dimexpr\linewidth-#1\relax}\vskip1pt}%
  {\vskip1pt\end{minipage}\end{lrbox}\fbox{\usebox{\wboxbox}}\par\nointerlineskip}

\newcommand{\Исполнитель}{Иванов Иван Иванович}
\newcommand{\Руководитель}{Петров Пётр Петрович}
\newcommand{\Тема}{Производственная практика, Преддипломная}
\newcommand{\ВидПрактики}{Производственная практика}
\newcommand{\ТипПрактики}{Производственная практика, Преддипломная}
\newcommand{\Место}{Департамент математики, механики и компьютерных наук ФГАОУ ВО УрФУ ИЕНиМ}
\newcommand{\Институт}{ИЕНиМ}
\newcommand{\Кафедра}{школы бакалавриата (школа)}
\newcommand{\Код}{10.05.01}
\newcommand{\Специальность}{Компьютерная безопасность}
\newcommand{\Программа}{10.05.01/22.01 - Математические методы защиты информации}
\newcommand{\Группа}{МЕН-123456}
\newcommand{\Год}{2026}
\newcommand{\Начало}{14.09.\Год{}г.}
\newcommand{\Конец}{20.12.\Год{}г.}
\newcommand{\СдачаС}{11.12.\Год{}г.}
\newcommand{\СдачаПо}{14.12.\Год{}г.}
\newcommand{\Замечания}{Нет предложений и замечаний}

\makeatletter
\def\w@short#1 #2#3 #4#5\relax{#1 #2.#4.}
\newcommand{\Кратко}[1]{\edef\w@fio{#1}%
  \ifx\w@fio\empty\else\expandafter\w@short\w@fio\relax\fi}
\makeatother

\begin{document}
\begin{wpar}[\wcenter\wsz{12}{13.80}]ОТЗЫВ\end{wpar}
\пусто[\wjustify\wsz{12}{13.80}]
\begin{wpar}[\wjustify\wsz{12}{20.70}]ФИО \Исполнитель,\end{wpar}
\begin{wpar}[\wjustify\wsz{12}{20.70}]Вид практики: \ВидПрактики\end{wpar}
\begin{wpar}[\wjustify\wsz{12}{20.70}]Тип практики: \ТипПрактики\end{wpar}
\begin{wpar}[\wjustify\wsz{12}{20.70}]Организация:\end{wpar}
\пусто[\wjustify\wsz{12}{20.70}]
\begin{wpar}[\wsz{12}{20.70}]с \Начало{} по \Конец\end{wpar}
\пусто[\wsz{12}{13.80}]
\begin{wpar}[\wjustify\wsz{12}{20.70}]\wstretch{Студент(ка) \Исполнитель{} группы \Группа{} института \Институт{} за время прохождения практики осуществил(а) следующие мероприятия}\wdash[\wul{\phantom{\_}}]{308}\end{wpar}
\begin{wpar}[\wsz{12}{20.70}]В период практики студент(ка)\break \wdash[\wul{\phantom{\_}}]{385}\end{wpar}
\begin{wpar}[\wcenter\wsz{12}{20.70}]\wsize{8}{(краткая характеристика уровня подготовки и отношения практиканта к работе)}\end{wpar}
\пусто[\wjustify\wsz{12}{20.70}]
\begin{wpar}[\wjustify\wrule{1.5pt}\wsz{12}{20.70}]\textbf{Ответить на вопросы}\end{wpar}
\begin{wpar}[\wsz{12}{13.80}]1.Студент на время практики~был трудоустроен?\wspaces{3}\wfield{(да,нет)}\wspaces{3}\wfield{(указать должность)}\end{wpar}
\пусто[\wsz{12}{13.80}]
\begin{wpar}[\wsz{12}{13.80}]2.Студенту предложено пройти следующую практику на предприятии (в организации)? \wfield{(да,нет)}\end{wpar}
\пусто[\wsz{12}{13.80}]
\begin{wpar}[\wsz{12}{13.80}]3. Студенту предложено трудоустройство после завершения обучения \wfield{(да,нет)}\end{wpar}
\пусто[\wsz{12}{13.80}]
\begin{wpar}[\wsz{12}{13.80}]4. Предложения и замечания от организации по теоретической и практической подготовке студентов.\wspaces{2}\wcap{\wulp{\Замечания}}{\wsub{(В свободной форме)}}\end{wpar}
\пусто[\wsz{12}{13.80}]
\пусто[\wsz{12}{13.80}]
\begin{wpar}[\wsz{12}{20.70}]Оценка за практику \hbox{\wul{\phantom{\_\_\_\_\_\_\_}}}\end{wpar}
\begin{wpar}[\wsz{12}{20.70}]«\_\_\_\_\_\_\_\_\_\_»\_\_\_\_\_\_\_\_\_\_\_\_\_\_\_\_\_\_\_\_\_\_ 20\_\_\_\_г.\end{wpar}
\begin{wpar}[\wsz{12}{20.70}]Руководитель организации\wspaces{61}(подпись, печать)\_\_\_\end{wpar}
\begin{wpar}[\wsz{12}{20.70}]\wtab{\wtab{Руководитель практики от организации\wspaces{22}}}\wspaces{1}(подпись)\_\_\_\_\_\_\_\_\_\_\end{wpar}
\пусто[\wspace{0pt}{10pt}\wsz{12}{15.87}]
\end{document}
